%% kapa: exam=1 round=26 year=2015 period=2 form=A subject=accounting subject_ko=회계학 seq=22 number=62
%% source: Q-Net 감정평가사 기출문제 (cst00309) artlSeq=5138001 file=extracted/1008/5138001/제26회_감정평가사 1차_기출/제26회_감정평가사 1차_기출/2교시_감평_A형.pdf
%% conversion: kordoc 4.24.51 -> tools/gen1.py (auto); review=proofread-sonnet; flags=none
\begin{question}{62}
(주)감평은 20×5년 6월 2일에 주식 100주를 주당 ￦2,500에 취득하고, 취득과 직접적으로 관련된 수수료 ￦15,000을 지급하였다. 위 주식의 각 연도말 주당 공정가치는 다음과 같다.
\begin{framedtext}
\centering\begin{tabular}{ccc}
\uline{20×5년 말} & \uline{20×6년 말} & \uline{20×7년 말} \\
￦2,800 & ￦2,450 & ￦2,700
\end{tabular}\par
\end{framedtext}
(주)감평이 20×8년 1월 2일에 동 주식 전부를 주당 ￦2,750에 처분하였다. (주)감평이 동 주식을 취득시점에서 단기매매금융자산으로 분류했을 경우와 매도가능금융자산으로 분류했을 경우로 구분하여 20×8년에 인식할 처분손익을 각각 계산하면 얼마인가?

\noindent\hspace*{2em}\makebox[12em][l]{\uline{단기매매금융자산}}\uline{매도가능금융자산}\par
\begin{choices}
  \item \makebox[12em][l]{손실 ￦25,000}손실 ￦25,000
  \item \makebox[12em][l]{손실 ￦5,000}이익 ￦25,000
  \item \makebox[12em][l]{이익 ￦5,000}이익 ￦10,000
  \item \makebox[12em][l]{이익 ￦5,000}이익 ￦25,000
  \item \makebox[12em][l]{이익 ￦5,000}손실 ￦10,000
\end{choices}
\end{question}
\answer{3}
